\documentclass{article}
\usepackage{fontspec}
\setmainfont{Libertinus Serif}
\usepackage[a4paper, margin=1in]{geometry}
\usepackage{amsmath, amsthm, amssymb, mathtools, enumitem, hyperref}
\usepackage[portuguese]{babel}

\begin{document}

\begin{flushleft}
    UFRGS - Instituto de Matemática e Estatística \\
    Departamento de Matemática Pura e Aplicada \\
    MAT 01326 - Matemática Discreta I \\
\end{flushleft}
\vspace{0.25cm}

\begin{center}
    {\LARGE Lista 6 - Coloração} \\
    \vspace{0.4cm}
    {\large \textbf{Soluções}}

\end{center}

\subsection*{Colorindo vértices}

\begin{enumerate}
    \item [\textbf{E1}] Observa que em grafos bipartidos completos da forma \(K_{1,n}\), a diferença \(\Delta(G) - \chi(G)\) pode ser feita arbitrariamente grande para um grafo \(G\). Para que \(\Delta(G)-\chi(G)=5\), basta construir o grafo \(K_{1,7}\), onde \(\Delta(G)=7\) e \(\chi(G)=2\). Para \(n\) qualquer, faz \(K_{1, n+2}\).
    \item [\textbf{E2}]
    \item [\textbf{E3}] Dado \(\chi(G)=3\), construímos um grafo tripartido completo \(G\). A contagem de arestas para tal grafo \(G\) com \(3n\) vértices é \(\frac{1}{2}(3n)(2n)=3n^2\). Para um grafo com \(3n+1\) vértices, a contagem muda para \(3n^2+2n\).
\end{enumerate}

\subsection*{Colorindo arestas}

\begin{enumerate}
    \item [\textbf{E5}]
    \item [\textbf{E6}]
    \item [\textbf{E7}]
    \item [\textbf{E8}]
\end{enumerate}

\end{document}